% Capítulo 5: Resultados y discusión.

Este capítulo presenta los resultados como un recorrido que va de la
base experimental al diagnóstico final. Se describe primero la
construcción y verificación de los datos y la caracterización de los
regímenes de cada formulación. Sobre esa base se evalúan los métodos
clásicos y, después, el método cuántico, antes de compararlos bajo un
protocolo común. Las dos últimas secciones separan el diagnóstico entre
fallo del criterio y fallo del optimizador, y examinan la consistencia
de los resultados junto con sus limitaciones.

\section{Construcción y verificación de la base experimental}
\label{sec:res-base}

El primer resultado del trabajo es la propia base experimental: antes
de comparar selectores se verificó que los datos, el preprocesado y
las particiones no introducen sesgos que contaminen la comparación.

\subsection{Análisis exploratorio de los datos crudos}

El análisis exploratorio no detectó incidencias de calidad básicas
(ningún conjunto presenta valores ausentes, filas duplicadas ni columnas
constantes) y marcó para exclusión los identificadores estructurales
(\texttt{id} en \textit{Breast Cancer} y \textit{Olive Oil},
\texttt{CustomerID} en \textit{Customer Churn}). Resuelta esa
verificación, el interés está en la presión dimensional de partida que
resume la Figura~\ref{fig:dimensionalidad}: \textit{Madelon} dispone de
solo 4 filas por variable predictora, frente a las más de 44\,000 de
\textit{Customer Churn}.

\begin{figure}[!htbp]
  \centering
  \includegraphics[width=\textwidth]{F01_banco_regimenes.pdf}
  \caption[Composición del banco experimental]{Composición del banco experimental. El plano sitúa cada
    formulación dataset--target por número de muestras y de predictoras
    en escala logarítmica; el panel central muestra el ratio de
    desbalance mayoría--minoría en entrenamiento, que descarta la
    \textit{accuracy} como métrica común; y el panel derecho resume el
    porcentaje de variables con señal univariante tras corrección FDR\@.
    Esta heterogeneidad justifica usar macro-F1 como métrica común.}
  \label{fig:dimensionalidad}
\end{figure}
\FloatBarrier

La asociación univariante con el target, evaluada con los contrastes
descritos en la sección~\ref{sec:mt-validacion}
(Mann--Whitney/Kruskal--Wallis para variables numéricas y $\chi^2$ para
categóricas), mostró perfiles muy distintos. En \textit{Breast Cancer}
y \textit{Olive Oil} abundan variables con tamaños de efecto altos,
varios de ellos casi perfectos, lo que obligó a una revisión semántica
de posibles proxies; en \textit{Olive Oil} esa revisión condujo a la
separación de problemas descrita en la
sección~\ref{sec:oliveoil}. En \textit{Customer Churn}, el gran tamaño
muestral hace que todas las variables resulten estadísticamente
significativas, por lo que la lectura se desplazó a la magnitud del
efecto. La corrección de Benjamini--Hochberg~\cite{benjamini1995}
cuantifica esos perfiles (Figura~\ref{fig:senal-fdr-efecto}): sobreviven
27 de 30 variables en \textit{Breast
Cancer} (efecto absoluto mediano 0{,}65), las 10 de 10 en \textit{Customer
Churn} y \textit{Olive Oil} (efecto 0{,}20 y 0{,}80 respectivamente), pero
solo \textbf{13 de las 500} de \textit{Madelon}, con un efecto mediano de
apenas 0{,}02. Esta última cifra es la primera evidencia de que la señal
de \textit{Madelon} apenas es univariante (reside en interacciones de
orden superior), un diagnóstico que anticipa por qué los criterios
basados en información mutua, clásicos y cuántico, encontrarán ahí su
límite. La redundancia, medida con correlaciones de Spearman, resultó
estructural en \textit{Breast Cancer} (pares con $|\rho| > 0{,}95$ entre
medidas de la misma familia) y en la composición química de
\textit{Olive Oil}. La Tabla~\ref{tab:senal} sintetiza estos perfiles
de señal y redundancia por conjunto.

\begin{table}[!htbp]
  \centering
  \caption[Señal univariante y redundancia por conjunto]{Señal
    univariante y redundancia de los datos crudos. ``Signif.\ FDR'' es el
    número de variables que conservan asociación con el target tras la
    corrección de Benjamini--Hochberg sobre las contrastadas; el efecto
    mediano es el tamaño de efecto absoluto (delta de Cliff en problemas
    binarios numéricos, V de Cramér en categóricas y épsilon cuadrado en
    multiclase); la última columna cuenta los pares de variables con
    correlación de Spearman $|\rho|\ge0{,}85$.}
  \label{tab:senal}
  \scriptsize
  \resizebox{\textwidth}{!}{%
  \begin{tabular}{lrrrrr}
    \toprule
    Dataset & Variables & Signif.\ sin corr. & Signif.\ FDR &
    Efecto mediano & Pares $|\rho|\ge0{,}85$ \\
    \midrule
    Breast Cancer Wisconsin & 30 & 27 & 27 & 0{,}65 & 29 \\
    Customer Churn & 10 & 10 & 10 & 0{,}23 & 0 \\
    Madelon & 500 & 38 & 13 & 0{,}02 & 12 \\
    Olive Oil & 10 & 10 & 10 & 0{,}80 & 2 \\
    \bottomrule
  \end{tabular}%
  }
\end{table}
\FloatBarrier

\begin{figure}[!htbp]
  \centering
  \includegraphics[width=\textwidth]{F02_senal_fdr_efecto.pdf}
  \caption[Señal supervisada univariante]{Señal supervisada univariante por dataset. El panel de FDR
    muestra cuántas variables mantienen asociación con el target tras
    corregir multiplicidad; el panel de efecto separa significancia de
    magnitud. \textit{Madelon} queda resaltado porque casi no presenta
    señal univariante, anticipando el límite del criterio QFS.}
  \label{fig:senal-fdr-efecto}
\end{figure}
\FloatBarrier

\subsection{Preprocesado y comprobación posterior al preprocesado}

El preprocesado fue deliberadamente conservador: exclusión de
identificadores, normalización de nombres y codificación determinista
del target, sin imputación (no había nulos), sin tratamiento de
valores atípicos (se caracterizaron como evidencia descriptiva) y sin
codificación ni escalado global de predictoras, que se pospusieron al
interior de los pipelines posteriores al particionado.

La comprobación posterior verificó que esta intervención mínima no alteró
el problema (Tabla~\ref{tab:postproc}). Las proporciones del target se
conservaron exactamente en los cuatro conjuntos (diferencia máxima de
proporciones nula, contrastes $\chi^2$ y de Fisher con $p=1{,}0$ y delta
de entropía nula); las variables numéricas comparables mantuvieron
mediana y dispersión (Kolmogorov--Smirnov, Wasserstein y PSI nulos en
los cuatro conjuntos); y el orden de asociación variable--target se
conservó casi de forma perfecta, con una correlación de Spearman entre
los \textit{rankings} antes y después del preprocesado de 0{,}988 en el
peor caso (\textit{Olive Oil}) y una variación media de información
mutua inferior a 0{,}02. La estructura de redundancia tampoco se modificó,
y \textit{Customer Churn} llegó al particionado con sus tres predictoras
categóricas sin codificar, como exige el control de \textit{leakage}.

\begin{table}[!htbp]
  \centering
  \caption[Invarianza del preprocesado]{Verificación de invarianza del
    preprocesado. Todas las métricas comparan el dato crudo con el
    procesado: la correlación de Spearman entre los \textit{rankings} de
    asociación variable--target debe ser cercana a $1$, y el resto de
    estadísticos de desplazamiento (KS, Wasserstein, PSI) y la variación
    de redundancia deben ser nulos si el preprocesado no alteró el
    problema.}
  \label{tab:postproc}
  \scriptsize
  \resizebox{\textwidth}{!}{%
  \begin{tabular}{lrrrr}
    \toprule
    Dataset & $\rho$ Spearman \textit{rankings} & $\Delta$ MI media &
    KS / Wass.\ / PSI máx & $\Delta$ pares redund. \\
    \midrule
    Breast Cancer Wisconsin & 0{,}9996 & 0{,}002 & 0{,}00 & 0 \\
    Customer Churn & 1{,}000 & 0{,}003 & 0{,}00 & 0 \\
    Madelon & 1{,}000 & 0{,}000 & 0{,}00 & 0 \\
    Olive Oil & 0{,}988 & 0{,}018 & 0{,}00 & 0 \\
    \bottomrule
  \end{tabular}%
  }
\end{table}
\FloatBarrier

\subsection{Particionado y verificación de las particiones}

El particionado generó particiones 70/15/15 estratificadas para las
cinco formulaciones operativas. Las comprobaciones exactas no detectaron
solapes de filas entre particiones ni pérdida de clases: las nueve
regiones de \texttt{olive\_oil\_9class} están presentes en
entrenamiento, validación y prueba.

Una partición solo es informativa si entrenamiento, validación y prueba
proceden de la misma distribución; de lo contrario, el rendimiento medido
en prueba no sería extrapolable. El \emph{drift} es justamente ese
desplazamiento de la distribución de una variable entre dos particiones.
Se cribó variable a variable comparando las muestras de entrenamiento y
prueba con los mismos estadísticos de desplazamiento del preprocesado
(Kolmogorov--Smirnov, Wasserstein y PSI), de modo que una variable se
marca cuando alguno supera su umbral. Un número alto de marcas no
invalida la partición, pero obliga a leerlo como una cautela de
representatividad, sobre todo en conjuntos pequeños, donde una única
región mal repartida basta para disparar el contraste. La
Tabla~\ref{tab:particiones} recoge esta verificación para las cinco
formulaciones.

\begin{table}[!htbp]
  \centering
  \caption[Verificación de las particiones]{Verificación de las cinco
    particiones. La diferencia máxima de proporciones del target y el
    PSI máximo miden el desplazamiento entre particiones; el AUC
    adversarial (clasificador entrenamiento--prueba con validación
    cruzada) se interpreta como no accionable cuando queda cerca de 0{,}5
    bajo el umbral operativo de la verificación; la última columna cuenta variables señaladas por la
    pantalla de \textit{leakage} (AUC o información mutua normalizada
    $\ge0{,}99$).}
  \label{tab:particiones}
  \scriptsize
  \resizebox{\textwidth}{!}{%
  \begin{tabular}{lrrrrr}
    \toprule
    Formulación & $\Delta$ máx target & Vars \textit{drift} & PSI máx &
    AUC adversarial & Vars \textit{leakage} \\
    \midrule
    Breast Cancer Wisconsin & 0{,}005 & 25 & 0{,}50 & 0{,}52 & 0 \\
    Customer Churn & 0{,}000 & 0 & 0{,}00 & 0{,}52 & 0 \\
    Madelon & 0{,}000 & 41 & 0{,}18 & 0{,}48 & 0 \\
    Olive Oil (3 clases) & 0{,}012 & 7 & 0{,}34 & 0{,}51 & 0 \\
    Olive Oil (9 clases) & 0{,}012 & 6 & 0{,}33 & 0{,}54 & 0 \\
    \bottomrule
  \end{tabular}%
  }
\end{table}
\FloatBarrier

El cribado univariante de \textit{drift} marcó para revisión algunas
variables (25 en \textit{Breast Cancer}, 41 en \textit{Madelon},
ninguna en \textit{Customer Churn}), con PSI máximo de 0{,}50; estas
marcas se interpretaron como cautelas de representatividad en conjuntos
pequeños, no como defectos del particionado. La validación adversarial
(Figura~\ref{fig:adversarial}) cierra la verificación: el AUC quedó entre
0{,}476 (\textit{Madelon}) y 0{,}535 (\texttt{olive\_oil\_9class}), en todos
los casos dentro del umbral operativo $|AUC-0{,}5|<0{,}05$ usado como señal
no accionable de distinguibilidad entre particiones. La pantalla de \textit{leakage} no encontró ninguna
variable con AUC o información mutua normalizada $\ge 0{,}99$ una vez
excluidos los proxies de \textit{Olive Oil}.

\begin{figure}[!htbp]
  \centering
  \includegraphics[width=\textwidth]{F03_base_confiable.pdf}
  \caption[Fiabilidad de la base experimental]{Fiabilidad de la base experimental. La figura cierra la
    puerta a explicaciones triviales antes de modelar: las particiones no
    son distinguibles de forma accionable (AUC adversarial cercano a
    $0{,}5$), el \textit{drift} se interpreta como cautela de
    representatividad y no como fuga, y la conservación de los rankings de
    información mutua muestra que el preprocesado no fabricó ni destruyó
    la señal que después reciben los selectores.}
  \label{fig:adversarial}
\end{figure}
\FloatBarrier

\subsection{Regímenes de los conjuntos de datos como predictores}
\label{sec:regimenes-predictor}

El análisis de las cuatro primeras fases no se interpreta solo como un
control de calidad. Sus métricas definen el régimen en el que se evaluará
QFS y, por tanto, anticipan qué componente del método puede convertirse
en cuello de botella. La Tabla~\ref{tab:regimenes} resume esa lectura
con cuatro ejes:
magnitud de señal, multicolinealidad, dimensionalidad intrínseca y
representatividad de las particiones. La multicolinealidad se mide con el
\emph{factor de inflación de la varianza} (VIF): para cada variable,
cuantifica cuánto aumenta la varianza de su coeficiente de regresión por
estar correlacionada con el resto. Un VIF igual a~$1$ indica ausencia de
colinealidad, y como umbrales habituales se toman valores por encima
de~$5$ (o de~$10$ en criterios más laxos) como señal de
multicolinealidad apreciable~\cite{james2013}.

\begin{table}[!htbp]
  \centering
  \caption[Regímenes del dato antes de QFS]{Regímenes del dato antes de
    ejecutar QFS\@. El efecto mediano resume la magnitud de la señal
    supervisada; el VIF máximo y el número de variables con VIF$\ge10$
    miden la multicolinealidad; los componentes principales para explicar
    el 80\,\% de la varianza aproximan la dimensionalidad intrínseca. Las
    dos formulaciones de \textit{Olive Oil} comparten predictoras.}
  \label{tab:regimenes}
  \scriptsize
  \resizebox{\textwidth}{!}{%
  \begin{tabular}{lrrrr}
    \toprule
    Dataset & Efecto mediano & VIF máx & Vars VIF$\ge10$ &
    Comp.\ PCA (80\,\%) \\
    \midrule
    Breast Cancer Wisconsin & 0{,}65 & 3806 & 23 & 5 \\
    Customer Churn & 0{,}23 & 1{,}1 & 0 & 6 \\
    Madelon & 0{,}02 & 117 & 20 & 295 \\
    Olive Oil & 0{,}80 & 326 & 5 & 3 \\
    \bottomrule
  \end{tabular}%
  }
\end{table}
\FloatBarrier

En \textit{Breast Cancer Wisconsin}, la señal supervisada vive en las
medidas morfológicas de la familia de tamaño y borde: entre las
relevancias $I_i$ más altas aparecen
\varname{concave_points_worst} ($0{,}415$),
\varname{radius_worst} ($0{,}404$) y
\varname{concave_points_mean} ($0{,}398$). El régimen no es difícil por
falta de señal, sino por redundancia: las familias
\texttt{radius}--\texttt{perimeter}--\texttt{area} alcanzan VIF máximo
de 3806 y 23 variables superan VIF 10. Esta estructura predice que un
selector de relevancia pura tenderá a escoger variables equivalentes,
mientras que mRMR y QFS deberían aportar si penalizan redundancia sin
perder la señal principal.

En \textit{Customer Churn}, la asociación con el target está concentrada
en pocas variables continuas o binarias interpretables:
\texttt{support\_calls} ($I_i=0{,}229$) domina, seguida de
\texttt{total\_spend} ($0{,}164$),
\texttt{contract\_length\_Monthly} ($0{,}130$),
\texttt{payment\_delay} ($0{,}092$) y \texttt{age} ($0{,}086$).
La redundancia es casi nula en las variables continuas (VIF máximo
1{,}1), pero la codificación \textit{one-hot} introduce grupos de dummies
mutuamente excluyentes, en especial \texttt{gender\_Female} frente a
\texttt{gender\_Male}. El régimen predice un criterio QUBO relativamente
fácil (hay poca redundancia que resolver) y un margen predictivo muy
pequeño porque el baseline roza la separación perfecta. Si aparece un
deterioro, debe buscarse en el readout u optimización, no en una falta
de señal del criterio.

\textit{Madelon} presenta el régimen opuesto. Solo 13 de 500 variables
superan FDR tras corrección, el efecto mediano es 0{,}02 y las relevancias
$I_i$ son casi planas: la mejor variable del handoff QFS,
\texttt{feat\_241}, alcanza $0{,}026$, y las siguientes
\texttt{feat\_475} y \texttt{feat\_338} quedan en $0{,}023$ y $0{,}020$.
Además, la PCA requiere 295 componentes para explicar el 80\,\% de la
varianza. La dificultad no es una multicolinealidad masiva que el término
de redundancia pueda depurar, sino un desierto de distractores donde la
señal útil depende de interacciones que una información mutua par a par
no ordena bien. Por ello, antes de ejecutar QFS ya se espera un fallo de
criterio si el QUBO hereda esas relevancias.

La Tabla~\ref{tab:relevancias} reúne, por conjunto, las variables de mayor
relevancia $I_i$ (información mutua con el target, calculada solo con
entrenamiento), que sostienen las lecturas anteriores y son las que el
método cuántico hereda como término lineal del QUBO\@. La
Figura~\ref{fig:ap-handoff} resume estas entradas de forma agregada: la
relevancia media $I_i$ y la redundancia media $R_{ij}$ por formulación,
tal como las recibe la fase cuántica.

\begin{figure}[!htbp]
  \centering
  \includegraphics[width=0.86\textwidth]{a6_handoff_ir.png}
  \caption[Handoff de relevancia y redundancia hacia QFS]{Entradas que
    recibe QFS: relevancia media $I_i$ con el target y redundancia media
    $R_{ij}$ entre pares de variables, por formulación. El orden por
    relevancia deja visible el contraste de regímenes: \textit{Olive Oil}
    y \textit{Breast Cancer} llegan con señal clara, mientras que
    \textit{Madelon} entra casi plano; por eso, si QFS falla allí, la
    primera sospecha debe recaer en el criterio que hereda esas
    relevancias, no en el gráfico ni en el clasificador posterior.}
  \label{fig:ap-handoff}
\end{figure}
\FloatBarrier

\begin{table}[!htbp]
  \centering
  \caption[Variables de mayor relevancia por conjunto]{Tres variables de
    mayor relevancia $I_i$ (información mutua con el target) por
    formulación, con su valor entre paréntesis. Es el vector que QFS recibe
    como término lineal del QUBO\@. La diferencia de escala entre conjuntos
    es informativa: en \textit{Customer Churn} \texttt{support\_calls}
    domina con holgura, mientras que en \textit{Madelon} ninguna variable
    supera $0{,}03$, lo que anticipa el límite del criterio.}
  \label{tab:relevancias}
  \small
  \begin{tabularx}{\textwidth}{@{}l>{\raggedright\arraybackslash}X>{\raggedright\arraybackslash}X>{\raggedright\arraybackslash}X@{}}
    \toprule
    Formulación & 1.\textsuperscript{a} & 2.\textsuperscript{a} & 3.\textsuperscript{a} \\
    \midrule
    Breast Cancer & \varname{concave_points_worst} (0,41) & \varname{radius_worst} (0,40) & \varname{concave_points_mean} (0,40) \\
    Customer Churn & \varname{support_calls} (0,23) & \varname{total_spend} (0,16) & \varname{contract_length_M.} (0,13) \\
    Madelon & \varname{feat_241} (0,026) & \varname{feat_475} (0,023) & \varname{feat_338} (0,020) \\
    Olive Oil (3 cl.) & \varname{other} (0,62) & \varname{linoleic} (0,41) & \varname{stearic} (0,36) \\
    Olive Oil (9 cl.) & \varname{linolenic} (0,87) & \varname{other} (0,78) & \varname{linoleic} (0,76) \\
    \bottomrule
  \end{tabularx}
\end{table}
\FloatBarrier

En \textit{Olive Oil}, las ocho variables son porcentajes de ácidos
grasos y forman un régimen composicional. Las relevancias se concentran
en \texttt{other}, \texttt{linolenic}, \texttt{linoleic} y
\texttt{stearic}; la redundancia es moderada y estructural porque los
ácidos forman partes de una composición. La formulación de tres clases
es casi separable. La de nueve clases mantiene las mismas variables,
pero el conjunto de prueba contiene solo 86 filas y varias regiones
químicamente próximas, de modo que el resultado debe leerse con
prudencia. Este régimen predice equivalencia o mejora no concluyente,
no un fallo atribuible del método.

En conjunto, la base experimental no define un único problema de
selección, sino cinco regímenes distintos: \textit{Breast Cancer} y las
dos formulaciones de \textit{Olive Oil} representan escenarios con señal
clara; \textit{Customer Churn}, un régimen de techo predictivo casi
saturado y baja redundancia; y \textit{Madelon}, un caso adverso para los
criterios marginales o par a par. Esta clasificación condiciona la
interpretación posterior de los métodos clásicos y del método cuántico:
no se espera el mismo comportamiento en todos los conjuntos, sino que el
protocolo permita identificar dónde y por qué aparece cada limitación.

\section{Evaluación de los métodos clásicos}

Esta sección evalúa los doce selectores clásicos sobre las cinco
formulaciones del banco. En primer lugar se caracteriza qué selecciona
cada método en el plano relevancia--redundancia; después se modela con
esos subconjuntos y se contrasta su rendimiento mediante intervalos de
confianza y pruebas pareadas. Las dos últimas subsecciones resumen la
comparación final en un mapa de evidencia y leen el espacio de métodos
clásicos como espejo del comportamiento esperable de QFS.

\subsection{Caracterización de la selección}

Los doce selectores del repertorio (sección~\ref{sec:clasicos}) se
ejecutaron sobre las cinco formulaciones, con tres semillas
independientes para los once que recorren una escalera de presupuestos y
una ejecución por dataset para Boruta, que confirma un conjunto de tamaño
fijo. Tres propiedades caracterizan la selección clásica más allá de los
rankings: la estabilidad entre semillas, la redundancia interna del
subconjunto y el coste de ajuste. La Tabla~\ref{tab:perfil} reúne las
tres por método y la Figura~\ref{fig:perfil-selectores} las representa.

\begin{table}[!htbp]
  \centering
  \caption[Perfil de los doce selectores]{Perfil de los doce selectores,
    promediado sobre las cinco formulaciones. El \emph{Jaccard medio} es
    el solape $|A\cap B|/|A\cup B|$ entre los subconjuntos obtenidos con
    pares de semillas distintas al presupuesto de referencia: vale $1$
    cuando la selección es idéntica entre semillas. La \emph{redundancia
    interna} es la correlación absoluta media entre las variables
    seleccionadas (menor es mejor). El \emph{coste} es el tiempo medio de
    ajuste por ejecución. Boruta no recorre semillas, por lo que no tiene
    Jaccard medio.}
  \label{tab:perfil}
  \small
  \begin{tabular}{lrrr}
    \toprule
    Método & Jaccard medio & Redundancia interna & Coste (s) \\
    \midrule
    \texttt{variance}            & 1{,}00 & 0{,}220 & 0{,}007 \\
    \texttt{f\_classif}          & 1{,}00 & 0{,}389 & 0{,}006 \\
    \texttt{mutual\_info}        & 0{,}99 & 0{,}388 & 0{,}102 \\
    \texttt{mutual\_correlation} & 1{,}00 & 0{,}365 & 0{,}116 \\
    \texttt{feature\_similarity} & 1{,}00 & 0{,}230 & 0{,}106 \\
    \texttt{mrmr\_approx}        & 0{,}95 & 0{,}167 & 0{,}265 \\
    \texttt{rrfs}                & 0{,}99 & 0{,}357 & 0{,}335 \\
    \texttt{l1\_logistic}        & 0{,}99 & 0{,}202 & 0{,}195 \\
    \texttt{linear\_svm}         & 0{,}99 & 0{,}276 & 0{,}011 \\
    \texttt{random\_forest}      & 0{,}89 & 0{,}384 & 0{,}560 \\
    \texttt{rfe}                 & 1{,}00 & 0{,}297 & 0{,}905 \\
    Boruta                       & ---  & 0{,}315 & 6{,}543 \\
    \bottomrule
  \end{tabular}
\end{table}
\FloatBarrier

\textbf{Estabilidad.} La selección es estable en las cinco
formulaciones: el Jaccard medio (solape entre los subconjuntos elegidos
con semillas distintas) es 1{,}0 en la mayoría de los métodos y solo baja
de forma apreciable en \textit{random forest} (0{,}89 de media, y 0{,}78 en
\textit{Madelon}). No aparece ningún régimen de inestabilidad, de modo
que los subconjuntos son reproducibles y el criterio que distingue a los
métodos no es la estabilidad, sino la redundancia interna.

Esta lectura admite una cautela. El Jaccard no corrige el solape
esperable por azar, que crece cuando el presupuesto es grande frente al
número de variables, como en \textit{Olive Oil} con sus ocho predictoras.
Por ello la estabilidad se contrastó también con el índice de
Kuncheva~\cite{kuncheva2007}, que sí corrige ese azar: coincide con el
Jaccard en los filtros y los métodos \emph{embedded}, pero rebaja la
estabilidad media de mRMR (de 0{,}95 a 0{,}62) y de la información mutua (de
0{,}99 a 0{,}64). Parte del solape entre semillas de esos dos criterios es,
por tanto, atribuible al azar en los conjuntos de pocas variables. Un
segundo índice independiente, la correlación de Spearman entre los
rankings de relevancia obtenidos con semillas distintas, apunta en la
misma dirección (0{,}75 y 0{,}74 de media para esos dos métodos), lo que
confirma que el efecto es real y no un artefacto de un único índice. Esta
menor estabilidad conviene tenerla presente al interpretar QFS, que se
apoya en ese mismo criterio de relevancia--redundancia; no afecta, en
cambio, a los veredictos de rendimiento, que se sostienen sobre p-valores
corregidos (FDR y Holm) e intervalos bootstrap, y no sobre la estabilidad
de la selección.

\begin{figure}[!htbp]
  \centering
  \includegraphics[width=\textwidth]{F04_perfil_selectores.pdf}
  \caption[Perfil de los doce selectores]{Perfil de los doce selectores. Los tres paneles usan el mismo
    orden para que la lectura sea comparativa: la línea discontinua marca
    la referencia de cada eje (redundancia media del espacio completo,
    coste de un segundo y estabilidad perfecta). La conclusión operativa
    es que \texttt{mrmr\_approx} ocupa el extremo de menor redundancia con
    coste moderado, mientras que Boruta paga mucha más computación por un
    perfil de redundancia menos favorable.}
  \label{fig:perfil-selectores}
\end{figure}
\FloatBarrier

Para que esa estabilidad no oculte la dificultad del problema, la Figura~\ref{fig:ap-organismo-seleccion} abre la caja negra de \textit{Madelon}: doce selectores recorren las
mismas 500 variables a lo largo de la escalera de presupuestos y tres
semillas. El consenso no se reparte de forma homogénea, sino que aparece
concentrado en un pequeño frente de variables mientras el resto del
espacio queda como distractor. Esta estructura explica por qué la
selección clásica puede mejorar el baseline y, a la vez, anticipa el
límite posterior de QFS: un criterio basado en información mutua por
pares puede ordenar bien parte del frente, pero no capturar todas las
interacciones que hacen útil a cada variable en este problema.

\begin{figure}[!htbp]
  \centering
  \includegraphics[width=\textwidth]{o1_organismo_seleccion.png}
  \caption[Organismo de selección en \textit{Madelon}]{Organismo de selección en \textit{Madelon}: frecuencia con
    que cada selector elige cada variable al recorrer presupuestos y
    semillas. La lámina conserva el frente de consenso que el cuerpo
    resume sin desplegar las 500 variables.}
  \label{fig:ap-organismo-seleccion}
\end{figure}
\FloatBarrier

El contraste por permutaciones del target aportó evidencia de cribado
compatible con señal en los cinco problemas: la mediana del p-valor
empírico de las variables seleccionadas alcanzó el mínimo posible
($\approx 0{,}048$) en todas las
formulaciones, incluida \textit{Madelon}. La mayoría de los métodos
aumenta la redundancia interna al concentrar variables correlacionadas,
mientras que \texttt{mrmr\_approx} es el único con un cambio
sistemáticamente nulo o negativo, confirmando el valor de penalizar la
redundancia de forma explícita.

\textbf{Coste y redundancia.} El perfil completo de los doce selectores
(coste, estabilidad y redundancia interna) se conserva en la Tabla~\ref{tab:perfil}; la lectura por ingrediente es la que importa
aquí. Los filtros univariantes son baratos (del orden de $0{,}007$\,s) pero
concentran redundancia ($\approx 0{,}39$ de correlación media), y los
\textit{wrappers} son dos a tres órdenes de magnitud más lentos (Boruta,
$6{,}5$\,s de media); frente a ellos, \texttt{mrmr\_approx} logra la menor
redundancia interna ($0{,}167$) a coste moderado. El método cuántico
persigue precisamente esa baja redundancia de \texttt{mrmr\_approx}
resolviendo el compromiso de forma global.

El cruce de las variables seleccionadas con la evidencia de las
comprobaciones previas no encontró ninguna variable con patrón de
identificador o etiqueta: las exclusiones del particionado llegaron
intactas a la selección. Los 168 subconjuntos (dataset~$\times$~método~$\times$~$k$)
mantienen columnas consistentes entre particiones y se acompañan del
grafo de características de la sección~\ref{sec:grafo} (vector $I_i$ y
matriz $R_{ij}$ por formulación, calculados solo con entrenamiento).

La evolución completa de la redundancia interna frente al presupuesto
$k$ queda como evidencia de control en la Figura~\ref{fig:ap-redundancia-k}; en el cuerpo basta con fijar la
conclusión que conecta con QFS: la penalización explícita de redundancia
es el ingrediente que distingue a \texttt{mrmr\_approx} de los filtros de
relevancia pura.

\begin{figure}[!htbp]
  \centering
  \includegraphics[width=\textwidth]{f5_coste_rendimiento.png}
  \caption[Redundancia interna frente a $k$]{Redundancia interna al crecer el presupuesto $k$ de
    variables. Las líneas grises muestran el contexto de selectores
    clásicos; \texttt{mrmr\_approx} y Boruta quedan resaltados como dos
    referencias de lectura. La comparación muestra por qué el criterio de
    redundancia importa para QFS: mRMR mantiene subconjuntos compactos y
    poco correlacionados, mientras que Boruta actúa como referencia
    \textit{wrapper} más costosa y no siempre más limpia en redundancia.}
  \label{fig:ap-redundancia-k}
\end{figure}
\FloatBarrier

\subsection{Modelado y evaluación estadística}

El modelado cruzó la parrilla de 260 experimentos (5 formulaciones
$\times$ 13 espacios de variables, el conjunto completo más los doce
selectores, $\times$ 4 modelos: regresión logística, SVM lineal,
\textit{random forest} y \textit{XGBoost}) y evaluó en prueba los 15
candidatos fijados con validación. La inclusión de \textit{XGBoost}
proporciona el anclaje de comparabilidad con la evaluación original del
método cuántico~\cite{orquin2026} y, al ser un clasificador de gradiente
robusto en alta dimensión, eleva de forma apreciable el rendimiento del
baseline en \textit{Madelon} (de 0{,}613 con modelos lineales o de bosque
a 0{,}813), lo que reduce, sin anularlo, el margen de mejora atribuible a la
selección. La Tabla~\ref{tab:comparacion} resume la comparación final por
dataset y la Figura~\ref{fig:test} muestra los intervalos de confianza
(el detalle por modelo y métrica de cada candidato se detalla en la Tabla~\ref{tab:candidatos}).

\begin{table}[!htbp]
  \centering
  \caption[Candidatos evaluados en prueba por dataset y modelo]{Candidatos
    evaluados en prueba: el mejor baseline y los dos mejores subconjuntos
    seleccionados en validación, con el modelo que los gana, el número de
    variables y las métricas en el conjunto de prueba. El AUC-ROC solo
    aplica a los dos problemas binarios; en los multiclase la métrica
    primaria es el macro-F1. En \textit{Customer Churn} las 10 predictoras
    crudas se materializan como 15 columnas tras la codificación
    \textit{one-hot}.}
  \label{tab:candidatos}
  \footnotesize
  \begin{tabular}{lllrrr}
    \toprule
    Dataset & Configuración & Modelo & vars & macro-F1 & AUC \\
    \midrule
    \multirow{3}{*}{Breast Cancer W.}
      & baseline      & SVM lineal & 30 & 0{,}937 & 0{,}976 \\
      & Boruta        & SVM lineal & 22 & 0{,}950 & 0{,}975 \\
      & \texttt{linear\_svm} & XGBoost & 10 & 0{,}950 & 0{,}970 \\
    \midrule
    \multirow{3}{*}{Customer Churn}
      & baseline      & XGBoost & 15 & 1{,}000 & 1{,}000 \\
      & $\ell_1$ logística & XGBoost & 10 & 0{,}999 & 1{,}000 \\
      & \texttt{mrmr\_approx} & XGBoost & 10 & 0{,}999 & 1{,}000 \\
    \midrule
    \multirow{3}{*}{Madelon}
      & baseline      & XGBoost & 500 & 0{,}813 & --- \\
      & Boruta        & XGBoost & 19 & 0{,}907 & --- \\
      & \texttt{rfe}  & XGBoost & 10 & 0{,}900 & --- \\
    \midrule
    \multirow{3}{*}{Olive Oil (3 cl.)}
      & baseline      & SVM lineal & 8 & 1{,}000 & --- \\
      & F de ANOVA    & SVM lineal & 5 & 1{,}000 & --- \\
      & $\ell_1$ logística & SVM lineal & 5 & 1{,}000 & --- \\
    \midrule
    \multirow{3}{*}{Olive Oil (9 cl.)}
      & baseline      & SVM lineal & 8 & 0{,}839 & --- \\
      & $\ell_1$ logística & XGBoost & 5 & 0{,}954 & --- \\
      & Boruta        & SVM lineal & 8 & 0{,}839 & --- \\
    \bottomrule
  \end{tabular}
\end{table}
\FloatBarrier

\begin{table}[!htbp]
  \centering
  \caption[Comparación final en test]{Comparación final en el conjunto de prueba: baseline con
    todas las variables frente al mejor subconjunto seleccionado. Se
    reporta el macro-F1, la diferencia pareada $\Delta$ con su p-valor
    corregido por comparaciones múltiples (Benjamini--Hochberg) y el
    veredicto. Un veredicto es de mejora o deterioro solo cuando el
    intervalo bootstrap del 95\,\% de $\Delta$ excluye el cero y la
    corrección lo confirma; en \textit{Customer Churn} la diferencia es
    significativa pero tan pequeña que se lee como empate práctico.}
  \label{tab:comparacion}
  \small
  \resizebox{\textwidth}{!}{%
  \begin{tabular}{lrrlrrrl}
    \toprule
    Dataset & \multicolumn{2}{c}{Baseline} &
    \multicolumn{3}{c}{Mejor selección} & $\Delta$ F1 ($p_{\text{FDR}}$) & Veredicto \\
    & vars & F1 & método & vars & F1 & & \\
    \midrule
    Breast Cancer W. & 30 & 0{,}937 & Boruta & 22 & 0{,}950 & $+0{,}013$ ($1{,}00$) & equivalente \\
    Customer Churn & 15 & 1{,}000 & $\ell_1$ logística & 10 & 0{,}999 & $-0{,}001$ ($0{,}001$) & empate práctico \\
    Madelon & 500 & 0{,}813 & Boruta & 19 & 0{,}907 & $+0{,}094$ ($0{,}001$) & mejora \\
    Olive Oil (3 clases) & 8 & 1{,}000 & F de ANOVA & 5 & 1{,}000 & $0{,}000$ ($1{,}00$) & equivalente \\
    Olive Oil (9 clases) & 8 & 0{,}839 & $\ell_1$ logística & 5 & 0{,}954 & $+0{,}116$ ($0{,}71$) & no concluyente \\
    \bottomrule
  \end{tabular}%
  }
\end{table}
\FloatBarrier

Los subconjuntos ganadores son interpretables variable a variable. En
\textit{Customer Churn}, la $\ell_1$ logística retiene
\varname{support_calls}, \varname{total_spend}, \varname{payment_delay},
\varname{age}, \varname{tenure}, \varname{last_interaction}, una dummy de
\varname{gender} y una de \varname{subscription_type}, junto a
\varname{contract_length_Monthly}: las continuas informativas más
representantes categóricos. En \textit{Olive Oil} de tres clases el
subconjunto de F de ANOVA es \{\varname{linoleic}, \varname{linolenic},
\varname{other}, \varname{palmitoleic}, \varname{stearic}\}, y en la
formulación de nueve clases la $\ell_1$ logística cambia
\varname{palmitoleic} y \varname{linoleic} por \varname{eicosenoic} y
\varname{oleic}. En \textit{Breast Cancer}, el conjunto confirmado por
Boruta (22 de 30) abarca las familias de \varname{radius}, \varname{area},
\varname{perimeter}, \varname{concavity} y \varname{concave_points} en
sus variantes \texttt{mean}, \texttt{se} y \texttt{worst}, y solo
descarta ocho medidas redundantes. En \textit{Madelon}, las 19 variables
de Boruta reúnen tanto la cabeza de la información mutua
(\varname{feat_241}, \varname{feat_475}) como las que el modelo realmente
explota (\varname{feat_336}, \varname{feat_105}, \varname{feat_153}), un
detalle que reaparece en el diagnóstico de la
sección~\ref{sec:diagnostico-qfs}. El roster completo por método se
conserva en la Figura~\ref{fig:ap-roster}.

\begin{figure}[!htbp]
  \centering
  \includegraphics[width=\textwidth]{a3_roster_completo.png}
  \caption[Consenso de selección por variable]{Consenso de la selección
    clásica: para cada variable, cuántos de los doce selectores la eligen
    en su presupuesto de referencia. Las barras de la izquierda (rojo)
    son variables que casi nadie elige (distractores); hacia la
    derecha, variables de consenso. En \textit{Madelon} domina el pico de
    variables elegidas por un solo método (el desierto de distractores),
    mientras que en los conjuntos con señal clara el consenso se reparte
    hacia la derecha.}
  \label{fig:ap-roster}
\end{figure}
\FloatBarrier

El veredicto por dataset aplica una regla de decisión con tres
condiciones simultáneas (diferencia positiva, intervalo bootstrap de la
diferencia pareada que excluye el cero y p-valor pareado corregido por
multiplicidad inferior a 0{,}05) y un umbral de efecto práctico de
$0{,}01$ en macro-F1: una diferencia estadísticamente significativa pero
de magnitud inferior a ese umbral se declara empate práctico. El valor
$0{,}01$ (un punto porcentual de macro-F1) se fija como mínima diferencia
operativamente relevante para preferir un conjunto de variables sobre
otro; por debajo de él, una diferencia puede ser significativa por el
tamaño muestral sin tener consecuencias prácticas, como ilustra
\textit{Customer Churn}.

\begin{figure}[!htbp]
  \centering
  \includegraphics[width=\textwidth]{F07_significancia_magnitud.pdf}
  \caption[Diferencia pareada en test]{Diferencia pareada de macro-F1 en test entre la mejor
    selección y el baseline, con intervalo \textit{bootstrap} del
    95\,\%. La línea discontinua marca el umbral práctico de 0{,}01, que
    separa la significancia estadística de la magnitud útil.}
  \label{fig:test}
\end{figure}
\FloatBarrier

El patrón es consistente (Tabla~\ref{tab:comparacion}): en cuatro de los
cinco problemas la selección iguala al baseline con menos variables, sin
que el intervalo de la diferencia pareada excluya el cero (en
\texttt{olive\_oil\_9class} la mejora aparente a 0{,}954 no es afirmable por
las 86 filas de prueba). \textit{Customer Churn} ilustra el umbral
práctico: el baseline saturado (1{,}000) supera al subconjunto (0{,}999) por
una diferencia significativa pero despreciable, que se declara empate.

\textit{Madelon} es el caso diferencial: el baseline con las 500
variables alcanza 0{,}813 y el conjunto confirmado por Boruta (19
variables) sube a 0{,}907. Es la única mejora estadísticamente
significativa del bloque tras corregir por comparaciones múltiples
($+0{,}094$, IC 95\,\% $[0{,}053, 0{,}140]$, p-valor corregido $0{,}001$ con FDR
y $0{,}005$ con Holm) y cuantifica el daño que los distractores infligen al
modelo completo, una vez que \textit{XGBoost} ya absorbe buena parte de
ese daño en el baseline. Los 15 candidatos superaron el contraste de
permutación de etiquetas con $p = 0{,}002$ (el mínimo alcanzable): ningún
resultado es compatible con un clasificador aleatorio. Las matrices de
confusión finales no muestran colapso de clases minoritarias; en
\texttt{olive\_oil\_9class} la confusión se concentra en dos regiones
que se reparten con sus vecinas.

La magnitud de la mejora de \textit{Madelon} depende del modelo. Antes
de fijar \textit{XGBoost} como baseline de comparación con QFS, los
modelos lineales o de bosque sufrían más los 480 distractores y la
selección producía mejoras cercanas a 0{,}28 de macro-F1. Con
\textit{XGBoost}, el baseline completo sube a 0{,}813 y la mejora de
Boruta se reduce a 0{,}094. Esta reducción no elimina el valor de
seleccionar, pero muestra que parte de la señal de interacción puede ser
recuperada por el modelo incluso cuando el criterio univariante no la
ordena. La comparación posterior con QFS debe leerse en ese contexto:
si el modelo usa variables que la información mutua de pares no sitúa
arriba, el límite será el criterio de selección, no la capacidad del
clasificador.

\textbf{Interpretabilidad.} Para cada candidato evaluado se calcularon
los valores SHAP (\textit{SHapley Additive exPlanations})~\cite{lundberg2017}
con el explicador nativo de su familia de modelo (\textit{TreeExplainer}
para bosques aleatorios y \textit{XGBoost}, \textit{LinearExplainer} para
regresión logística y SVM lineal), conservando la matriz completa por
instancia (y por clase en los problemas multiclase), no solo la
importancia media (Figura~\ref{fig:shap-variables}). Esta
evidencia confirma que las variables que sostienen cada modelo coinciden
con las que la selección retuvo y con las señaladas por la asociación
univariante del análisis exploratorio, cerrando el círculo entre
relevancia estadística, selección y contribución al modelo. En las formulaciones de
\textit{Olive Oil} el desglose por clase muestra que distintas regiones
se apoyan en ácidos grasos distintos, lo que explica la confusión
residual entre regiones vecinas. La lectura por instancia, además, es la
base de interpretabilidad sobre la que se comparará el subconjunto del
método cuántico.

La excepción que da sentido al diagnóstico posterior es
\textit{Madelon}. En ese dataset, los mayores valores SHAP incluyen
\texttt{feat\_336}, \texttt{feat\_105} y \texttt{feat\_153}, variables
que el vector $I_i$ de QFS no coloca en la cabeza con la misma claridad
que \texttt{feat\_241}. Esta discrepancia entre lo que el modelo usa y
lo que la relevancia por pares ordena es la huella variable a variable
del fallo de criterio que se observa después en el óptimo exacto del
QUBO.

\begin{figure}[!htbp]
  \centering
  \includegraphics[width=\textwidth]{F06_shap_beeswarm_instancias.pdf}
  \caption[Variables que sostienen los modelos]{Variables que sostienen los modelos elegidos por validación:
    distribución de valores SHAP por instancia (\textit{beeswarm}) en
    \textit{Breast Cancer}, \textit{Customer Churn} y \textit{Madelon}.
    Cada punto es una instancia del candidato seleccionado y el color
    codifica el valor de la variable, de bajo (azul) a alto (rojo).}
  \label{fig:shap-variables}
\end{figure}
\FloatBarrier

\subsection{Comparación final y mapa de evidencia}

La aportación de la selección clásica no es uniforme entre regímenes: es
máxima cuando abundan los distractores (\textit{Madelon}), neutra cuando
el problema ya es casi separable (\textit{Breast Cancer} y \textit{Olive
Oil} de tres clases) y apenas mueve el rendimiento cuando el baseline
está saturado (\textit{Customer Churn}). Esta lectura por regímenes es la
que condiciona después la comparación con el método cuántico.

La comparación final consolida la lectura conjunta del bloque sobre las
cuatro coordenadas que articulan la tesis. La
Figura~\ref{fig:perfil-selectores}, la Tabla~\ref{tab:perfil} y el
detalle de la Figura~\ref{fig:ap-redundancia-k} muestran que la estabilidad entre semillas es
uniforme y alta en todos los métodos, de modo que no discrimina entre
ellos; el verdadero diferenciador es la redundancia interna y su relación
con el rendimiento, a coste contenido salvo en los \textit{wrappers}.
Este contraste es el que conecta con el método cuántico, cuyo valor
diferencial es la optimización global del compromiso
relevancia--redundancia.

Como cierre operativo, el bloque clásico fija la referencia que el método
cuántico debe igualar o superar, y lo hace en las mismas cuatro
coordenadas, no solo en rendimiento: presupuesto de $k=10$ variables en
los tres conjuntos grandes y $k=5$ en las formulaciones de
\textit{Olive Oil}, macro-F1 de prueba de la Tabla~\ref{tab:comparacion},
redundancia interna comparable a la de \texttt{mrmr\_approx},
estabilidad de la selección y los mismos contrastes estadísticos de este
bloque. En cuatro problemas el objetivo es igualar el rendimiento con
igual o menor redundancia interna y un subconjunto más compacto; en
\textit{Madelon}, igualar o superar la mejora del mejor selector
clásico. Además de la
referencia de rendimiento, el bloque fija las entradas del método
cuántico (el grafo de características de la sección~\ref{sec:grafo}),
de modo que la fase cuántica no necesita recalcular ninguna cantidad
informacional y cualquier diferencia de resultados será atribuible al
método de optimización y no a las entradas.

\section{Evaluación del método cuántico}
\label{sec:res-cuantico}

Conviene fijar antes la notación que se repite en esta sección y en la
comparación posterior: \emph{QFS-NA} designa el subconjunto que devuelve
la simulación analógica de átomos neutros; el \emph{oráculo} (o
\emph{QFS-oráculo}) designa el óptimo exacto de la misma función QUBO al
presupuesto fijado, que sirve de control para separar el criterio del
optimizador; $\alpha$ controla en ese QUBO exacto el balance entre
relevancia y redundancia, y con él la cardinalidad del subconjunto; y
$\beta$ es el grado de libertad operativo del embebido que separa pares
de variables relevantes.

El método QFS se ejecutó conforme al diseño de la sección~\ref{sec:qfs}.
Las formulaciones que caben en la envolvente operativa entraron
directamente: \textit{Olive Oil}, con sus 8 predictores finales tras
excluir la etiqueta alternativa y el proxy \varname{palmitic}
(\varname{palmitoleic}, \varname{stearic}, \varname{oleic},
\varname{linoleic}, \varname{linolenic}, \varname{arachidic},
\varname{eicosenoic} y \varname{other}), y \textit{Customer Churn},
cuyas 10 predictoras crudas se materializan como 15 columnas tras la
codificación \textit{one-hot} de \varname{gender}, \varname{subscription_type}
y \varname{contract_length} (sección~\ref{sec:regimenes-predictor}). Las
que exceden la envolvente recibieron una preselección clásica previa con
\texttt{mrmr\_approx}, ajustada solo con entrenamiento (\textit{Breast
Cancer}, $30\rightarrow20$; \textit{Madelon}, $500\rightarrow20$),
declarada como parte del método híbrido. Por cada formulación se barrió el parámetro
$\beta$ en $\{0, 0{,}1, \ldots, 1\}$ y se eligió su valor con la
partición de validación usando \textit{random forest} como evaluador
fijo. La Figura~\ref{fig:qfs-beta} no resume solo un ajuste de
hiperparámetros: separa los dos mandos del método, el control exacto de
la cardinalidad mediante $\alpha$ en el QUBO y la reordenación operativa
inducida por $\beta$ en las densidades de Rydberg.

El embebido por escalado multidimensional se ejecutó con el protocolo
verificado de 100 inicializaciones independientes y 100 iteraciones
externas, guardando las posiciones y el error de embebido de cada
dataset, con la razón de distancia mínima de la referencia, $1/\sqrt{2}$ del
radio de bloqueo. La geometría se mide explícitamente, pero no se
utiliza como explicación antes de contrastarla con el error MDS.

\subsection{El mecanismo de QFS sobre los conjuntos de datos}
\label{sec:mecanismo-qfs}

Antes de leer rendimientos conviene ver \emph{cómo} el método convierte
cada conjunto en una decisión. La cadena es la misma en las cinco
formulaciones: el grafo de características (relevancia $I_i$ y redundancia
$R_{ij}$) se materializa en una matriz de átomos cuyas distancias salen
del embebido MDS; la evolución adiabática deja a cada átomo con una
\emph{densidad de Rydberg} media, su propensión a quedar excitado y, por
tanto, seleccionado; y el corte top-$k$ sobre esa densidad fija el
subconjunto. La Figura~\ref{fig:organismo-cuantico} hace visible esa
cadena: cada panel es una formulación, cada punto un átomo en su posición
real del embebido, el color es la densidad de Rydberg y el anillo marca
lo que el corte selecciona.

\begin{figure}[!htbp]
  \centering
  \includegraphics[width=\textwidth]{qfs_organismo_cuantico.pdf}
  \caption[Organismo cuántico: densidad de Rydberg y selección por
    formulación]{Organismo cuántico. Para cada formulación, los átomos
    (variables) se sitúan en su posición del embebido MDS, coloreados por
    su densidad de Rydberg media; el anillo negro marca los seleccionados
    por el corte top-$k$. El método elige, en esencia, los átomos de mayor
    densidad.}
  \label{fig:organismo-cuantico}
\end{figure}
\FloatBarrier

La lectura confirma el mecanismo y anticipa los diagnósticos posteriores.
El patrón general es nítido: los átomos seleccionados son casi siempre los
de mayor densidad de Rydberg, de modo que la densidad \emph{es} el ranking
operativo del método. Sobre ese patrón, cada régimen deja su firma. En
\textit{Customer Churn}, parte del presupuesto se va a las \textit{dummies}
de \texttt{subscription\_type} y de \texttt{gender}, que alcanzan densidad
alta y desplazan a continuas como \texttt{tenure}: es la huella visual del
déficit de optimizador que se cuantifica en la
sección~\ref{sec:diagnostico-qfs}. En \textit{Madelon}, el corte sí captura
átomos de densidad alta (\texttt{feat\_241}, \texttt{feat\_205},
\texttt{feat\_346}), lo que muestra que el optimizador hace su trabajo; el
problema, como se verá, es que esa densidad hereda un criterio $I_i$ casi
plano y no apunta a las variables que el modelo necesita. \textit{Breast
Cancer} concentra la densidad en la familia de tamaño y borde celular que
ya dominaba la relevancia, y las dos formulaciones de \textit{Olive Oil}
muestran un mapa pequeño y composicional donde la selección recae sobre los
ácidos de mayor densidad (\texttt{other}, \texttt{linolenic},
\texttt{arachidic}).

\subsection{Parámetros de control de QFS y comparación con el óptimo exacto}
\label{sec:mandos-control-qfs}

El método tiene dos mandos que conviene no confundir
(Figura~\ref{fig:qfs-beta}): el parámetro $\alpha$, que en el QUBO exacto
traza una escalera discreta de cardinalidades, y el parámetro
$\beta$ de la ejecución analógica, que reordena las densidades de Rydberg.

\begin{figure}[!htbp]
  \centering
  \includegraphics[width=\textwidth]{F08_mandos_qfs_alpha_beta.pdf}
  \caption[Mandos de QFS: $\alpha$ y $\beta$]{Mandos de QFS\@. El barrido de $\alpha$ corresponde al óptimo
    exacto del QUBO y se usa como control clásico del criterio, no como
    salida del simulador analógico. El barrido de $\beta$ pertenece a la
    ejecución QFS-NA y modifica la geometría de distancias, el rendimiento
    de validación y las densidades de Rydberg de las variables.}
  \label{fig:qfs-beta}
\end{figure}
\FloatBarrier

Junto a la ejecución operativa del simulador analógico ($\alpha$ fijo y
corte top-$k$ sobre la densidad de Rydberg), se resolvió el óptimo exacto
de la función de coste $Q(x;\alpha)$ por enumeración dentro de la
envolvente, recorriendo $\alpha$ según el procedimiento de
Mücke~\cite{mucke2023}. Esta referencia exacta es el mecanismo que
convierte un posible resultado nulo en evidencia interpretable: permite
separar si el simulador optimiza bien su criterio y si ese criterio, aun
optimizado, produce subconjuntos predictivos. La
Tabla~\ref{tab:qfs-control} resume ese control de calidad.

\begin{table}[!htbp]
  \centering
  \caption[Control de calidad del optimizador analógico]{Control de
    calidad del optimizador analógico frente al control exacto del QUBO a
    presupuesto fijo, al valor de $\beta$ elegido por validación. La distancia de Hamming
    cuenta en cuántas variables difiere el subconjunto del simulador
    (QFS-NA) del control exacto \texttt{mucke\_k} sobre el mismo espacio de
    la envolvente; el $\Delta$ coste evalúa ambos subconjuntos en
    $\alpha=0{,}5$ para medir coherencia con la ejecución operativa. Como el
    control \texttt{mucke\_k} se obtiene con restricción de cardinalidad y
    un $\alpha$ propio, no es necesariamente el óptimo libre de
    $\alpha=0{,}5$; por eso $\Delta$ puede ser negativo.}
  \label{tab:qfs-control}
  \small
  \begin{tabular}{lrr}
    \toprule
    Dataset & Hamming & $\Delta$ coste \\
    \midrule
    Breast Cancer W. & 6 & $-4{,}399$ \\
    Customer Churn & 6 & $+1{,}194$ \\
    Madelon & 2 & $-0{,}008$ \\
    Olive Oil (3 clases) & 2 & $-0{,}165$ \\
    Olive Oil (9 clases) & 2 & $+0{,}164$ \\
    \bottomrule
  \end{tabular}
\end{table}
\FloatBarrier

\begin{figure}[!htbp]
  \centering
  \includegraphics[width=\textwidth]{F09_atribucion_criterio_optimizador.pdf}
  \caption[Diagnóstico criterio--optimizador]{Diagnóstico criterio--optimizador en unidades de macro-F1. El
    eje vertical mide la pérdida atribuible al criterio
    relevancia--redundancia (baseline frente al control exacto del QUBO
    a presupuesto fijo);
    el eje horizontal mide la pérdida atribuible al optimizador o lectura
    analógica (control exacto frente a QFS-NA). Así, \textit{Madelon} y
    \textit{Customer Churn} no son dos fallos equivalentes: uno queda
    limitado por el criterio y el otro por la ejecución simulada.}
  \label{fig:criterio-optimizador}
\end{figure}
\FloatBarrier

La lectura del control es diagnóstica, pero no se interpreta de forma
aislada. En \textit{Madelon}, el coste de QFS-NA no explica el deterioro
principal: a $\alpha=0{,}5$ la diferencia de coste es prácticamente nula y
ligeramente favorable ($-0{,}008$), mientras que el óptimo exacto que se evalúa más
abajo queda lejos del baseline. El problema no es alcanzar el mínimo de
la función de coste, sino que esa función, construida con $I_i$ casi
planos, no representa la señal que el modelo necesita. En
\textit{Customer Churn}, el coste sí queda por encima del óptimo
($+1{,}194$), pero el deterioro predictivo asociado es pequeño. Esta
distinción evita leer la evaluación cuántica como dos fallos simétricos:
uno es un límite del criterio; el otro, un déficit leve del readout u
optimización analógica.

\section{Comparación entre enfoques}
\label{sec:res-comparacion}

Los subconjuntos del método cuántico se evaluaron con el protocolo de
modelado idéntico al del bloque clásico (los cuatro modelos, macro-F1,
intervalos \textit{bootstrap} y contrastes pareados), reutilizando las predicciones
por fila ya archivadas para comparar contra el baseline sin reentrenarlo.
La Tabla~\ref{tab:qfs-comparacion} recoge, por formulación, la mejor
configuración operativa del QFS y la configuración del control exacto
(oráculo a presupuesto fijo) frente al baseline.

\begin{table}[!htbp]
  \centering
  \caption[Comparación QFS--baseline en test]{Comparación QFS--baseline en prueba (macro-F1). QFS-NA es la
    ejecución operativa del simulador analógico; QFS-oráculo es el control
    exacto de $Q(x;\alpha)$ a presupuesto fijo. El veredicto aplica la
    misma regla de decisión que el bloque clásico.}
  \label{tab:qfs-comparacion}
  \small
  \resizebox{\textwidth}{!}{%
  \begin{tabular}{lrrrrl}
    \toprule
    Dataset & Baseline & QFS-NA & QFS-oráculo & $\Delta$ QFS-NA ($p_{\text{FDR}}$) & Veredicto (QFS-NA) \\
    \midrule
    Breast Cancer W. & 0{,}937 & 0{,}937 & 0{,}937 & $+0{,}000$ ($1{,}00$) & equivalente \\
    Customer Churn & 1{,}000 & 0{,}987 & 0{,}999 & $-0{,}013$ ($0{,}001$) & deterioro \\
    Madelon & 0{,}813 & 0{,}647 & 0{,}643 & $-0{,}166$ ($0{,}001$) & deterioro \\
    Olive Oil (3 clases) & 1{,}000 & 1{,}000 & 1{,}000 & $+0{,}000$ ($1{,}00$) & equivalente \\
    Olive Oil (9 clases) & 0{,}839 & 0{,}854 & 0{,}906 & $+0{,}016$ ($1{,}00$) & no concluyente \\
    \bottomrule
  \end{tabular}%
  }
\end{table}
\FloatBarrier

El veredicto se reporta con el mejor modelo en validación, pero conviene
comprobar que no depende de esa elección. La Tabla~\ref{tab:qfs-robustez}
resume el rango de macro-F1 de QFS-NA entre los cuatro modelos: la
conclusión es estable salvo en \textit{Customer Churn}, donde los modelos
de árbol sostienen el rendimiento (0{,}99) y los lineales caen a 0{,}89, de
modo que el deterioro leve es, en parte, una interacción
modelo--subconjunto. La rejilla completa de los cuatro modelos por
configuración se conserva en la Tabla~\ref{tab:ap-qfs-modelos}.

\begin{table}[!htbp]
  \centering
  \caption[Rendimiento de QFS-NA por modelo]{Macro-F1 de prueba del
    subconjunto operativo QFS-NA con los cuatro clasificadores, por
    formulación. Resume la robustez del veredicto frente a la elección de
    modelo (Tabla~\ref{tab:qfs-robustez} del cuerpo).}
  \label{tab:ap-qfs-modelos}
  \small
  \begin{tabular}{lcccc}
    \toprule
    Dataset & Reg.\ logística & SVM lineal & Random forest & XGBoost \\
    \midrule
    Breast Cancer W. & 0{,}938 & 0{,}938 & 0{,}937 & 0{,}950 \\
    Customer Churn & 0{,}889 & 0{,}889 & 0{,}985 & 0{,}987 \\
    Madelon & 0{,}640 & 0{,}647 & 0{,}610 & 0{,}597 \\
    Olive Oil (3 clases) & 0{,}989 & 1{,}000 & 1{,}000 & 1{,}000 \\
    Olive Oil (9 clases) & 0{,}969 & 0{,}874 & 0{,}854 & 0{,}951 \\
    \bottomrule
  \end{tabular}
\end{table}
\FloatBarrier

\begin{table}[!htbp]
  \centering
  \caption[Robustez del veredicto de QFS-NA entre modelos]{Robustez del
    veredicto de QFS-NA frente al modelo. Rango de macro-F1 de prueba
    entre los cuatro clasificadores y modelo que lo maximiza.}
  \label{tab:qfs-robustez}
  \small
  \begin{tabular}{llcc}
    \toprule
    Dataset & Mejor modelo & Rango macro-F1 & Veredicto estable \\
    \midrule
    Breast Cancer W. & XGBoost & 0{,}937--0{,}950 & sí \\
    Customer Churn & XGBoost & 0{,}889--0{,}987 & no (lineales caen) \\
    Madelon & SVM lineal & 0{,}597--0{,}647 & sí (deterioro) \\
    Olive Oil (3 clases) & SVM/RF/XGB & 0{,}989--1{,}000 & sí \\
    Olive Oil (9 clases) & reg.\ logística & 0{,}854--0{,}969 & sí (inconcluso) \\
    \bottomrule
  \end{tabular}
\end{table}
\FloatBarrier

\begin{figure}[!htbp]
  \centering
  \includegraphics[width=\textwidth]{F10_comparacion_final_clasico_qfs.pdf}
  \caption[Comparación final clásico--QFS en test]{Comparación final en test entre baseline, mejor referencia
    clásica, QFS-NA y QFS-oráculo. QFS-NA es la ejecución operativa del
    simulador analógico; QFS-oráculo es el control exacto clásico del
    QUBO a presupuesto fijo. Las barras horizontales muestran el
    intervalo bootstrap del
    delta frente al baseline, y la separación entre QFS-NA y oráculo
    enlaza con la atribución de la Figura~\ref{fig:criterio-optimizador}.}
  \label{fig:qfs-vs-clasico}
\end{figure}
\FloatBarrier

\clearpage

Para comprobar si el veredicto anterior era solo un punto aislado del
presupuesto, se realizó un re-modelado dirigido sobre la escalera de
$k$: baseline, \texttt{mrmr\_approx} y QFS-NA ordenado por densidad de
Rydberg, manteniendo el \textit{random forest} y la métrica de macro-F1.
La trayectoria completa de validación se conserva en la Figura~\ref{fig:ap-rendimiento-k}.

\begin{figure}[!htbp]
  \centering
  \includegraphics[width=\textwidth]{ev6_rendimiento_vs_k.png}
  \caption[Rendimiento frente al tamaño del subconjunto]{Macro-F1 de validación frente al tamaño del subconjunto en el
    re-modelado dirigido. Se comparan baseline, mRMR y QFS-NA ordenado
    por densidad de Rydberg para comprobar si el presupuesto elegido es
    un punto aislado o parte de una trayectoria estable.}
  \label{fig:ap-rendimiento-k}
\end{figure}
\FloatBarrier

La comparación final (Figura~\ref{fig:qfs-vs-clasico}) admite una lectura
por regímenes, no una clasificación lineal de ganadores. En los problemas
con señal clara y espacio reducido (\textit{Breast Cancer}, \textit{Olive
Oil} de tres clases), el QFS iguala al baseline con subconjuntos
compactos. En \textit{Olive Oil} de nueve clases, la
configuración del oráculo sube hasta 0{,}906, pero el reducido conjunto de
prueba impide afirmar la mejora, en coherencia con la prudencia ya
aplicada al bloque clásico. Los deterioros informativos son
\textit{Madelon} y \textit{Customer Churn}, pero no tienen la misma
magnitud ni el mismo mecanismo.

El presupuesto $k$ y el recorrido de $\alpha$ permiten leer qué variable
entra antes en el criterio exacto (Figura~\ref{fig:ap-alpha}). En \textit{Customer Churn}, la
escalera de $\alpha$ comienza con \texttt{support\_calls}; en
$\alpha=0{,}5$ ya incorpora \texttt{payment\_delay},
\texttt{total\_spend} y \texttt{contract\_length\_Monthly}; y en
$\alpha=0{,}7$ añade \texttt{age}, \texttt{tenure} y una sola dummy de
\texttt{gender}. El orden del oráculo es coherente con el régimen:
prioriza continuas informativas y solo después introduce representantes
categóricos. QFS-NA, en cambio, retiene el grupo completo de
\texttt{subscription\_type} pero ya incorpora \texttt{payment\_delay} y
\texttt{total\_spend}; deja fuera \texttt{tenure} y
\texttt{last\_interaction}. El problema no es el
presupuesto $k$ en abstracto, sino qué cupos se consumen.

\begin{figure}[!htbp]
  \centering
  \includegraphics[width=\textwidth]{f10_b4_escalera_alpha.png}
  \caption[Escalera de $\alpha$ del óptimo exacto]{Escalera de $\alpha$ del óptimo exacto. El mecanismo traza una
    trayectoria discreta de cardinalidad que puede saltar presupuestos;
    en \textit{Madelon}, además, la entrada $I_i$ es casi plana y el
    orden de mérito deja de ser informativo.}
  \label{fig:ap-alpha}
\end{figure}
\FloatBarrier

En \textit{Madelon}, la misma escalera muestra otro mecanismo. El óptimo
exacto selecciona solo \texttt{feat\_241} desde $\alpha=0{,}1$ hasta
$\alpha=0{,}6$, pasa a dos variables en $\alpha=0{,}7$--$0{,}8$, cuatro en
$\alpha=0{,}9$ y salta a veinte en $\alpha=1{,}0$. La función de coste traza
una trayectoria de cardinalidades con saltos; para comparar a $k=10$, el
control exacto usa la restricción explícita de presupuesto. Aun así,
recibe un vector $I_i$ casi degenerado: \texttt{feat\_241}=0{,}026,
\texttt{feat\_475}=0{,}023 y \texttt{feat\_338}=0{,}020. Aumentar $k$ no
resuelve este caso si el criterio no distingue las variables que
interaccionan. El fallo de \textit{Madelon} es anterior al optimizador.

El parámetro $\beta$ sí reordena las densidades de Rydberg, pero no
explica el deterioro de \textit{Customer Churn}. La
Figura~\ref{fig:ap-beta}
muestra que los valores elegidos por validación son $\beta=0{,}1$ en
\textit{Breast Cancer}, $\beta=0{,}9$ en \textit{Customer Churn},
$\beta=0{,}1$ en \textit{Madelon}, $\beta=0{,}0$ en \textit{Olive Oil} de
tres clases y $\beta=0{,}2$ en la formulación de nueve clases. Todas las
configuraciones usan la razón de distancia $1/\sqrt{2}$ de la
referencia.

\begin{figure}[!htbp]
  \centering
  \includegraphics[width=\textwidth]{f10_b5_beta_geometria.png}
  \caption[Densidad de Rydberg frente a $\beta$]{Densidad de Rydberg media de las variables seleccionadas
    frente al resto al variar $\beta$, por dataset. La figura conserva
    el detalle operativo del mando geométrico de QFS-NA.}
  \label{fig:ap-beta}
\end{figure}
\FloatBarrier

La comprobación decisiva es el embebido MDS verificado, con 100
inicializaciones independientes (Tabla~\ref{tab:embedding}). El error
medio de embebido de \textit{Customer Churn} (0{,}234) es prácticamente
igual al de \textit{Breast Cancer} (0{,}225), que no se deteriora, y menor
que el de \textit{Madelon} (0{,}247). Por tanto, los datos verificados no
respaldan la hipótesis levantada durante el análisis (que el
\textit{one-hot} generase una geometría especialmente frustrada y
explicase el deterioro): una diferencia de embebido de una centésima
frente a un conjunto que iguala al baseline no basta para atribuir la
caída al MDS\@. El
\textit{one-hot} sigue siendo una divergencia de preprocesado relevante y
explica por qué existen grupos de dummies competidores, pero no se
sostiene como causa geométrica demostrada. El déficit de \textit{Customer
Churn} debe leerse como optimizador o \textit{readout} en un régimen de
baja redundancia y techo saturado, no como fallo de MDS\@. El embebido
atómico verificado queda documentado en la Figura~\ref{fig:ap-atomos-mds}.


\begin{table}[!htbp]
  \centering
  \caption[Error de embebido MDS por formulación]{Error medio de embebido
    por escalado multidimensional al $\beta$ elegido por validación, con
    100 inicializaciones independientes y razón de distancia
    $1/\sqrt{2}$ del radio de bloqueo. Un error mayor indica una
    geometría más difícil de realizar en el plano. El error de
    \textit{Customer Churn} es similar al de \textit{Breast Cancer}
    (que no se deteriora) y menor que el de \textit{Madelon}, lo que
    no respalda la geometría como explicación principal de su deterioro.}
  \label{tab:embedding}
  \small
  \begin{tabular}{lrr}
    \toprule
    Formulación & $\beta$ elegido & Error MDS medio \\
    \midrule
    Breast Cancer Wisconsin & 0{,}1 & 0{,}225 \\
    Customer Churn & 0{,}9 & 0{,}234 \\
    Madelon & 0{,}1 & 0{,}247 \\
    Olive Oil (3 clases) & 0{,}0 & 0{,}125 \\
    Olive Oil (9 clases) & 0{,}2 & 0{,}127 \\
    \bottomrule
  \end{tabular}
\end{table}
\FloatBarrier

\subsection{QFS en el espacio de selectores y en las cuatro coordenadas}
\label{sec:espacio-metodos-qfs}

La comparación no se agota en el macro-F1. La referencia clásica era un
sistema de coordenadas (relevancia, redundancia, estabilidad y coste), y
QFS debe leerse en ellas. En relevancia y redundancia, QFS se comporta
como una variante de la familia mRMR cuando el régimen lo permite: en
\textit{Olive Oil} de tres clases, QFS-NA y mRMR coinciden por completo
(Jaccard $1{,}00$, Tabla~\ref{tab:qfs-solape}) sobre los ácidos de alta
relevancia. En trazabilidad, el subconjunto QFS-NA comparte entradas
fijas con el bloque clásico, aunque su reproducibilidad completa no se
estima porque no se realizaron re-ejecuciones completas de la fase
cuántica; y en coste, la vía simulada es
órdenes de magnitud más lenta que cualquier filtro (Figura~\ref{fig:ap-coste-cuantico}), un coste que hoy solo se justifica
como exploración de viabilidad.


La Tabla~\ref{tab:qfs-solape} sitúa el subconjunto operativo frente a su
propio óptimo exacto y a mRMR (el solape frente a los doce métodos se
recoge en la Figura~\ref{fig:ap-jaccard-12}). La lectura
por variable es la informativa: en \textit{Madelon}, QFS-NA y el oráculo
comparten nueve de diez variables (la planitud de $I_i$ impide fijar un
subconjunto único, de ahí el bajo solape con mRMR); en \textit{Customer
Churn}, QFS-NA coincide con el oráculo solo en siete de diez: retiene
\varname{payment_delay} y \varname{total_spend} pero gasta cupo en
\textit{dummies} de \texttt{subscription\_type} y deja fuera
\texttt{tenure} y \varname{last_interaction}, que sí entran en el óptimo
exacto. Ese desplazamiento anticipa, sin probar todavía nada geométrico,
el déficit leve de optimizador del diagnóstico siguiente.


\begin{table}[!htbp]
  \centering
  \caption[Solape del subconjunto QFS-NA con el oráculo y con mRMR]{Solape
    (Jaccard) del subconjunto operativo QFS-NA con su óptimo exacto
    (oráculo) y con mRMR, al presupuesto de referencia. La última columna
    cuenta las variables compartidas con el oráculo.}
  \label{tab:qfs-solape}
  \small
  \begin{tabular}{lccc}
    \toprule
    Dataset & $J$(NA, oráculo) & $J$(NA, mRMR) & Comparten c/ oráculo \\
    \midrule
    Breast Cancer W. & 0{,}54 & 0{,}25 & 7 / 10 \\
    Customer Churn & 0{,}54 & 0{,}54 & 7 / 10 \\
    Madelon & 0{,}82 & 0{,}25 & 9 / 10 \\
    Olive Oil (3 clases) & 0{,}67 & 1{,}00 & 4 / 5 \\
    Olive Oil (9 clases) & 0{,}67 & 0{,}43 & 4 / 5 \\
    \bottomrule
  \end{tabular}
\end{table}
\FloatBarrier

\section{Diagnóstico criterio--optimizador}
\label{sec:diagnostico-qfs}

El plano de la Figura~\ref{fig:criterio-optimizador} resume la tesis
del bloque cuántico, y la Tabla~\ref{tab:atribucion} lo descompone en
cifras: para cada formulación separa la pérdida atribuible al criterio
(baseline frente al control exacto del QUBO a presupuesto fijo) de la
atribuible al optimizador o \textit{readout} (control exacto frente a
QFS-NA).

\begin{table}[!htbp]
  \centering
  \caption[Atribución criterio--optimizador]{Descomposición del
    rendimiento de QFS en macro-F1 de prueba. $\Delta$ criterio
    $=$ baseline $-$ oráculo mide cuánto pierde el subconjunto por usar el
    criterio relevancia--redundancia, aun resuelto de forma exacta;
    $\Delta$ optimizador $=$ oráculo $-$ QFS-NA mide cuánto pierde además
    la ejecución analógica respecto a ese óptimo. Un valor negativo indica
    mejora en lugar de pérdida. La separación entre ambas columnas es lo
    que distingue un fallo de criterio (\textit{Madelon}) de un déficit de
    optimizador (\textit{Customer Churn}).}
  \label{tab:atribucion}
  \small
  \begin{tabular}{lrrrrr}
    \toprule
    Dataset & Baseline & Oráculo & QFS-NA &
    $\Delta$ criterio & $\Delta$ optimizador \\
    \midrule
    Breast Cancer W. & 0{,}937 & 0{,}937 & 0{,}937 & $0{,}000$ & $0{,}000$ \\
    Customer Churn & 0{,}9998 & 0{,}9991 & 0{,}987 & $0{,}0007$ & $0{,}012$ \\
    Madelon & 0{,}813 & 0{,}643 & 0{,}647 & $0{,}170$ & $-0{,}003$ \\
    Olive Oil (3 clases) & 1{,}000 & 1{,}000 & 1{,}000 & $0{,}000$ & $0{,}000$ \\
    Olive Oil (9 clases) & 0{,}839 & 0{,}906 & 0{,}854 & $-0{,}067$ & $0{,}052$ \\
    \bottomrule
  \end{tabular}
\end{table}
\FloatBarrier

En \textit{Madelon}, el baseline con todas las
variables alcanza 0{,}813, el control exacto del QUBO baja a 0{,}643 y
QFS-NA queda en 0{,}647. La pérdida total es 0{,}166 de macro-F1, casi toda
atribuible al criterio (0{,}170): el optimizador no añade pérdida, pues
QFS-NA iguala al control exacto del QUBO a presupuesto fijo (0{,}647 frente
a 0{,}643). La explicación a nivel de variable coincide con SHAP: el modelo usa
\texttt{feat\_336}, \texttt{feat\_105} y \texttt{feat\_153}, entre
otras, mientras que el criterio MI del QFS rankea primero
\texttt{feat\_241} y apenas separa el resto. El modelo encuentra señal
en variables que la relevancia univariante no ordena bien. Por eso el
fallo de \textit{Madelon} no se corrige con más MDS ni con otro
$\beta$: es un límite de la formulación basada en información mutua por
pares.

En \textit{Customer Churn}, la pérdida total es mucho menor. El baseline
alcanza 0{,}9998, el oráculo 0{,}9991 y QFS-NA 0{,}987: el criterio solo
explica 0{,}0007 de pérdida, mientras que el optimizador explica 0{,}012.
El modelo sigue apoyándose en \texttt{support\_calls}, \texttt{age},
\texttt{total\_spend} y \texttt{contract\_length\_Monthly}, y todas
ellas permanecen en la selección QFS-NA, lo que explica que el deterioro
sea leve. La parte que falta se concentra en cupos gastados en el grupo
de dummies de \texttt{subscription\_type} y en la ausencia de
\texttt{tenure} y \texttt{last\_interaction}. Esta lectura descarta la simetría de dos
fallos iguales: \textit{Madelon} es un fallo grande de criterio;
\textit{Customer Churn} es un déficit leve de optimizador sin evidencia
geométrica.

El organismo cuántico (Figura~\ref{fig:organismo-cuantico}) hace visible
esa distinción a nivel de átomo. En \textit{Customer Churn}, las
\textit{dummies} de \texttt{subscription\_type} alcanzan densidades de
Rydberg altas y son seleccionadas, mientras que \texttt{tenure} y
\texttt{last\_interaction}, que el oráculo sí incorpora, quedan con
densidad baja: el fallo no es que el criterio ignore esas continuas, sino
que la lectura analógica de densidades las deja por debajo del corte. Es,
literalmente, un problema de \textit{readout}. En \textit{Madelon}, en
cambio, el corte selecciona átomos de densidad genuinamente alta
(\texttt{feat\_241}, \texttt{feat\_205}): el optimizador acierta con su
objetivo, pero ese objetivo, construido sobre un $I_i$ casi plano, no
coincide con las variables que el modelo necesita. La misma figura, dos
diagnósticos opuestos: en Churn la densidad no refleja bien al criterio;
en Madelon lo refleja con fidelidad, y el criterio es el que falla.

El resto de datasets cierra el marco. \textit{Breast Cancer} iguala al
baseline y las variables con mayor contribución SHAP
(\varname{texture_worst}, \varname{radius_worst},
\varname{area_worst}, \varname{concave_points_worst}) están dentro de
la familia que QFS retiene. \textit{Olive Oil} de tres clases alcanza
1{,}000 con ácidos de alta relevancia (\texttt{other},
\texttt{linolenic}, \texttt{linoleic}, \texttt{stearic}), y la versión
de nueve clases queda inconclusa por el tamaño de prueba. En este último
caso, macro-F1 es la métrica primaria porque el AUC binario no aplica a
la formulación multiclase; en los binarios se registró AUC solo como
puente con la referencia. \textit{Customer Churn} también cae en AUC
(0{,}983 frente a 1{,}000), lo que indica que la pérdida no es un artefacto
de macro-F1, aunque sigue siendo pequeña en un problema casi saturado.

\section{Consistencia, pruebas y discusión}
\label{sec:consistencia-discusion}

La atribución anterior depende de que se hayan descartado explicaciones
alternativas. La cadena de tests lo hace de forma explícita
(Tabla~\ref{tab:cadena-tests}): FDR y tamaño
de efecto separan señal real de significancia por tamaño muestral; la
comprobación posterior al preprocesado verifica que el preprocesado no altera
rankings ni target; la verificación de particiones descarta solapes,
\textit{leakage} y particiones adversarialmente distinguibles; la
selección se compara con 20 permutaciones del target, con p-valor mínimo
$1/21 \approx 0{,}048$; y el modelado se evalúa con 400 remuestreos
bootstrap, 2\,000 permutaciones de signos y 500 permutaciones de
etiquetas, cuya resolución mínima es $1/501 \approx 0{,}002$. Estos
eslabones no son accesorios: autorizan a atribuir \textit{Madelon} al
criterio y \textit{Customer Churn} al optimizador en lugar de al azar, al
\textit{split} o a una fuga de información.

\begin{table}[!htbp]
  \centering
  \caption[Cadena de controles estadísticos]{Cadena de controles estadísticos que acota las explicaciones alternativas antes del
    diagnóstico. Cada fila indica qué confusor descarta el control, su
    parámetro o resolución mínima de p-valor, y la evidencia donde se
    materializa en la memoria.}
  \label{tab:cadena-tests}
  \small
  \begin{tabular}{p{0.27\textwidth}p{0.34\textwidth}p{0.30\textwidth}}
    \toprule
    Control & Qué controla & Parámetro / evidencia \\
    \midrule
    Corrección FDR (Benjamini--Hochberg) & Falsos positivos por
    multiplicidad & $q=0{,}05$; Tabla~\ref{tab:senal} \\
    Tamaño de efecto & Significancia sin magnitud relevante & Cliff /
    $\varepsilon^2$ / Cramér; Tabla~\ref{tab:senal} \\
    Invarianza del preprocesado & Sesgo introducido al procesar &
    Spearman, KS, PSI; Tabla~\ref{tab:postproc} \\
    Verificación de particiones & \textit{Leakage}, \textit{drift}, split
    sesgado & AUC adversarial; Tabla~\ref{tab:particiones} \\
    Permutación del target (selección) & Señal espuria del selector &
    20 perm., $p_{\min}=1/21\approx0{,}048$ \\
    \textit{Bootstrap} de prueba & Incertidumbre de las métricas &
    400 remuestreos; Figura~\ref{fig:test} \\
    Permutación de signos pareada & Mejora por azar fila a fila &
    2\,000 perm.; Tabla~\ref{tab:comparacion} \\
    Permutación de etiquetas & Compatibilidad con clasificador aleatorio &
    500 perm., $p_{\min}=1/501\approx0{,}002$ \\
    Óptimo exacto del QUBO & Confusión criterio vs.\ optimizador &
    Enumeración; Tablas~\ref{tab:qfs-control} y~\ref{tab:atribucion} \\
    \bottomrule
  \end{tabular}
\end{table}
\FloatBarrier

La consistencia se midió donde el diseño lo permite. En selección, tres
semillas muestran que los filtros y métodos basados en información mutua
son muy estables entre semillas: su Jaccard medio (Tabla~\ref{tab:perfil}) es 1{,}00 para varios filtros al presupuesto de referencia. Ese solape bruto, sin embargo, sobreestima la estabilidad de mRMR y de la información mutua, cuyo índice de Kuncheva corregido por azar es sensiblemente menor. La variación aparece en métodos
estocásticos o en empates de variables casi equivalentes: \textit{random
forest} baja a 0{,}88 en \textit{Madelon} y 0{,}78 en
\textit{Olive Oil} de nueve clases; mRMR en \textit{Customer Churn}
queda en 0{,}88. La lectura variable a variable es benigna: en
\textit{Madelon} cambian representantes redundantes como
\texttt{feat\_28} y \texttt{feat\_48}; en \textit{Customer Churn} pueden
alternarse dummies de \texttt{subscription\_type}; en \textit{Olive Oil}
se intercambian ácidos composicionalmente próximos como
\texttt{palmitoleic} y \texttt{stearic}. Cambia el representante, no la
familia de información.

La Figura~\ref{fig:ap-consistencia} muestra esta estabilidad de forma
granular y contrasta el solape de Jaccard con el índice de Kuncheva: la
frialdad de mRMR e información mutua en el segundo confirma, conjunto a
conjunto, la cautela ya señalada sobre su estabilidad. Esta misma
fragilidad de la información mutua reaparece en el bloque cuántico:
\textit{Madelon}, donde su estabilidad corregida es mínima, es también el
conjunto en el que el criterio QUBO basado en información mutua fracasa en
rendimiento (comparación con el óptimo exacto, sección posterior), de modo
que estabilidad y rendimiento apuntan a la misma limitación del criterio
ante señal de orden superior.

\begin{figure}[!htbp]
  \centering
  \includegraphics[width=\textwidth]{f10_b10_consistencia.png}
  \caption[Estabilidad de la selección clásica: Jaccard frente a Kuncheva]{Estabilidad
    de la selección clásica entre semillas. Los dos paneles apilados
    separan la lectura bruta de la lectura corregida: el Jaccard sugiere
    estabilidad alta en casi todos los selectores, pero Kuncheva descuenta
    el solape esperable por azar y enfría mRMR e información mutua (en
    \textit{Madelon}, hasta $0{,}16$ y $0{,}19$). Esta diferencia es
    metodológicamente importante porque QFS hereda el mismo tipo de
    criterio relevancia--redundancia: donde la estabilidad corregida cae,
    la evidencia posterior debe interpretarse como posible límite del
    criterio, no como simple variación aleatoria.}
  \label{fig:ap-consistencia}
\end{figure}
\FloatBarrier

La consistencia de modelado y QFS se interpreta con prudencia. El
modelado no es un experimento multi-semilla; su robustez procede de
bootstrap, permutaciones pareadas y permutaciones de etiquetas sobre el
conjunto de prueba. QFS se ejecuta una vez por configuración de
$\beta$, pero cada embebido usa 100 inicializaciones MDS y 10\,000
mediciones del simulador. Una consistencia fuerte por múltiples
re-ejecuciones completas de modelado y QFS sería un experimento nuevo,
no una propiedad que deba inferirse de los artefactos actuales. Esta
limitación se declara porque forma parte del camino metodológico:
primero se aseguró la trazabilidad del criterio, del embebido y de las
predicciones; después se interpretó lo que esos controles autorizan.

En conjunto, los resultados sostienen una conclusión más precisa que
``QFS empata o pierde''. El bloque clásico predice el régimen; el espejo
de doce métodos sitúa a QFS en el compromiso relevancia--redundancia; la
escalera de $\alpha$ muestra si el criterio tiene un orden de mérito; el
embebido MDS no aporta evidencia suficiente para explicar geométricamente
la caída de \textit{Customer Churn}; SHAP muestra si las variables seleccionadas son las que usa el
modelo; y la cadena de tests permite atribuir con cautela. El resultado
final es que \textit{Madelon} queda limitado por el criterio, porque la
MI por pares no ve la señal de interacción; \textit{Customer Churn}
queda limitado de forma leve por el optimizador o readout en un régimen
de baja redundancia y techo saturado; \textit{Breast Cancer} y
\textit{Olive Oil} de tres clases no muestran fallo atribuible; y
\textit{Olive Oil} de nueve clases sigue abierto por tamaño muestral.

Esta lectura tiene, además, una consecuencia práctica que conviene
anticipar: allí donde QFS iguala a la selección clásica, el criterio
cuántico es competitivo, y esa competitividad es precisamente lo que
justifica explorar su ventaja computacional sobre hardware de átomos
neutros (seleccionar características sobre tantas variables como cúbits,
de forma global y en un tiempo de ejecución muy corto). Esta
proyección, y el estado actual de las plataformas de átomos neutros, se
desarrollan en el capítulo de conclusiones.
